\documentclass[10pt]{article}

\usepackage[utf8]{inputenc}

\usepackage[T1]{fontenc}

\usepackage{amsmath}

\usepackage{amsfonts}

\usepackage{amssymb}

\usepackage[version=4]{mhchem}

\usepackage{stmaryrd}

\usepackage{mathrsfs}

\usepackage{bbold}

\usepackage{graphicx}

\usepackage[export]{adjustbox}

\graphicspath{ {./images/} }



\begin{document}

ЧАСТОТНАЯ ТЕОРЕМА - теорема, формулирующая условия разрешимости у $\mathbf{p}$ а в н е н и й $Л$ у $\mathrm{p}$ b е



\[

P^{*} H+H P+h h^{*}=G, \quad H q-h \varkappa=q,

\]



где $P, G=G^{*}, q, g, x-$ заданные матрицы размеров $n \times n, n \times n, n \times m, n \times m, m \times m$ соответственно, $H=H^{*}$, $h$ - искомые матрицы размера $n \times n$ и $n \times m$. Уравнения Лурье имеют две другие әквивалентные формы: при $\operatorname{det} x \neq 0$



\[

H Q_{0} H+\left(P_{0}^{*} H+I I P_{0}\right)+G_{0}=0,

\]



где $Q_{0}=Q_{0}^{*} \geqslant 0, G_{0}=G_{0}^{*}$, и в общем случае



$2 \operatorname{Re} x^{*} H(\boldsymbol{P} x+q \xi)=\zeta(x, \xi)-\left|h^{*} x-x \xi\right|^{2} \quad(\forall x, \xi),(3)$ где $\mathscr{G}(x, \xi)-$ заданная әрмитова форма векторов $x \in \mathbb{C}^{n}$, $\xi \in \mathbb{C}^{m}:$



\[

\zeta(x, \xi)=x^{*} G x+2 \operatorname{Re}\left(x^{*} g \xi\right)+\xi^{*} \Gamma \xi .

\]



При этом $\Gamma=x^{*} x \geqslant 0, G_{0}=g \Gamma^{-1} g^{*}-G, P_{0}=P-g \Gamma g^{*}$, $Q_{0}=q \Gamma^{-1} q^{*}$.



Если пара $\{P, q\}$ управляема, то уравнения Лурье сводятся к случаю, когда



\[

P=\operatorname{diag}\left[\lambda_{1}, \ldots, \lambda_{n}\right], \quad \lambda_{j}+\lambda_{h} \neq 0 .

\]



При $m=1$ и когда все матрицы действительны, в скалярной записи уравнения Лурье приобретают вид



\[

\sum_{k=1}^{n} q_{k} \frac{h_{j} h_{k}}{\lambda_{j}+\lambda_{k}}-h_{j} V \bar{\Gamma}=\gamma_{j}, \quad j=1, \ldots, n,

\]



здесь $h=\left[h_{1}, \ldots, h_{n}\right]$ - искомый вектор.



Пусть пара $\{P, q\}$ с т а б и ли и и р у е ма, т. е. существует $r$ такое, что $R=P+q r^{*}-$ матрица Гурвица.



ч а с т о т в а я т е о р е м а утверждает: для разрешимости уравнений Лурье необходимо и достаточно, чтобы



\[

\mathscr{G}\left[(i \omega I-P)^{-1} q \xi, \xi\right] \geqslant 0

\]



для всех $\xi \in \mathbb{C}^{m}, \omega \in \mathbb{R}^{1}, \operatorname{det}\|i \omega I-P\| \neq 0$ (I- единичіная матрица). '.т. также формулирует процедуру определения матриц $I, h$ и утверждает, что при



$\operatorname{det} \Gamma \neq 0, \operatorname{det}\|i \omega I-P\| \neq 0, \quad \mathcal{G}\left[\|i \omega I-P\|^{-1} q \xi, \xi\right]>0$



\[

(\forall \xi \neq 0, \forall \omega)

\]



существуют такие (единственные) матрицы $H, h$, что кроме (3) выполнено: $P+q x^{-1} h^{*}$ есть матрица Гурвица (см. [3]).



Уравнение Лурье в форме (2) иногда наз. также мат тичным алгеб ра ическим ура вв е в и е м Р и к к а т и. Ч. т. иснолььзуется при решевии задач абсолютной устойчивости $[2,4,5]$, управления и адаптации (см., напр., [6]).



Лит.: [1] JI у р в е А. И., Некоторые нелинейные задачи теории автоматического регулирования, М.- Л., 1951; [2] II о пI о в В. М., Гиперустойчивость автоматических систем, пер. с рум., М., 1970; [3] Я к у б о н и ч В. А:, "Сиб. матем. ж.",



\includegraphics[max width=\textwidth, center]{2023_07_09_a146c4260c5161602b6ag-1(1)}

н о в Г. А., Я к у б о в и ч В. А., Устойчивость не.инсйных систем с неединственным состоянием равновесия, М., 1978; [5] Методы исследования нелинейных систем автоматического управления, М., 1975; [6] Ф о м и н В. Н., Ф р а д к о в А. JI., Я к у б о в ич В. А., Адаптивнос управјение динамическими



\includegraphics[max width=\textwidth, center]{2023_07_09_a146c4260c5161602b6ag-1(3)}

терполяционная квадратурная формула с равными коәффициентами:



\[

\int_{-1}^{1} f(x) d x \cong C^{\prime} \sum_{k=1}^{N} \int\left(x_{k}\right)

\]



Весовая функция равна 1 , промежуток интегрирования конечен и считается совпадающим с $[-1,1]$. Число параметров, определяющих квадратурную формулу (\textit{), равно $N+1$ ( $N$ узлов и значение коэффициента $C)$. Параметры определяются требованием, чтобы квадратурная формула (}) была точна для всех многочленов степени не выше $N$ или, что то же самое, для одночленов 1 , $x, x^{2}, \ldots, x^{N}$. Параметр $C$ находится из условия, что квадратурная формула точна для $f(x)=1$, и равен $2 / N$. Узлы $x_{1}, \ldots, x_{N}$ оказываются действительными лишь піри $N=1(1) 7$ и $N=9$. При $N=1(1) 7$ узлы вычислил П. Ј. Чебышев. При $N \geqslant 10$ среди узлов Ч. к. ф. всегда имеются комплексные (см. [1]). Алгебраич. степень точности Ч. к. ф. равна $N$ при $N$ нечетном и равна $N+1$ при $N$ четном. Формула (*) иредложена ІІ. ЈI. Чебышевым в 1873.



JІит.: [1] К р ы л о в В. И., Іриближенпое вычислсние И. ПТ. Мысовских.



ЧЕБЫIIЕВА МЕТОД - метољ иолучения класса итерациолыих алгоритмов нахождения однократного дейіствитенного корня уравнения



\[

f(x)=0,

\]



где $f(x)$ - достаточно гладкая функция.



В основе метода лежит формальное представление обратной к $f(x)$ фувкции $x=F(y)$ по формуле Тейлора. Если $\alpha$ - достаточно точное приближение для корня $x$ уравнения (1), $c_{0}=\beta=f(\alpha), f^{\prime}(\alpha) \neq 0$, то



\[

x=F(y)=\alpha+\sum_{n \geqslant 1} d_{n}(y-\beta)^{n},

\]



где коэффициенты $d_{n}$ рекуррентно определяются из соотношения $x \equiv F(f(x))$ через коэффициенты Тейлора $c_{n}$ функции $f(x)=\sum_{n \geqslant 0} c_{n}(x-\alpha)^{n}$. Полагая в (2) $y=0$, полутают соотвошение



\[

\begin{gathered}

x=\alpha-\frac{f(\alpha)}{f^{\prime}(\alpha)}-\left(\frac{f(\alpha)}{f^{\prime}(\alpha)}\right)^{2} \frac{f^{\prime \prime}(\alpha)}{2 f^{\prime}(\alpha)}- \\

-\left(\frac{f(\alpha)}{f^{\prime}(\alpha)}\right)^{3}\left(\frac{1}{2}\left(\frac{f^{\prime \prime}(\alpha)}{f^{\prime}(\alpha)}\right)^{2}-\frac{f^{\prime \prime \prime}(\alpha)}{6 f^{\prime}(\alpha)}\right)-

\end{gathered}

\]



$-\left(\frac{f(\alpha)}{f^{\prime}(\alpha)}\right)^{4}\left(\frac{5}{8}\left(\frac{f^{\prime \prime}(\alpha)}{f^{\prime}(\alpha)}\right)^{3}-\frac{5 f^{\prime \prime}(\alpha) f^{\prime \prime \prime}(\alpha)}{12\left(f^{\prime}(\alpha)\right)^{2}}+\frac{f^{1 \mathrm{~V}}(\alpha)}{24 f^{\prime}(\alpha)}\right)-\ldots$



Несколько членов справа в (3) дают формулы итерационного алгоритма; так при двух членах получается Ньютона меюоод, а при трех членах получается итерационный метод вида



$x_{n+1}=x_{n}-\frac{f\left(x_{n}\right)}{f^{\prime}\left(x_{n}\right)}-\left(\frac{f\left(x_{n}\right)}{f^{\prime}\left(x_{n}\right)}\right)^{2} \frac{f^{\prime \prime}\left(x_{n}\right)}{2 f^{\prime}\left(x_{n}\right)}, n=0,1, \ldots$



С ростом числа учитываемых в (3) членов возрастает скорость сходимости $x_{n}$ к $x$ (см. [2]). Метод может быть распространен на функциональные уравнения (см. [3]).



\begin{center}

\includegraphics[max width=\textwidth]{2023_07_09_a146c4260c5161602b6ag-1(4)}

\end{center}



\begin{center}

\includegraphics[max width=\textwidth]{2023_07_09_a146c4260c5161602b6ag-1(5)}

\end{center}



\begin{center}

\includegraphics[max width=\textwidth]{2023_07_09_a146c4260c5161602b6ag-1(2)}

\end{center}



\begin{center}

\includegraphics[max width=\textwidth]{2023_07_09_a146c4260c5161602b6ag-1}

\end{center}



ЧЕБЫІІЕВА МНОГО'ЛЕНЫ ІІ д а - многочлены, ортогональные на отрезке $[-1,1]$ с весовой функциеї



\[

h_{1}(x)=\frac{1}{\sqrt{1-x^{2}}}, \quad x \in(\cdots-1,1) .

\]



Для стандартизованных Ч. м. сираведливы формула



\[

T_{n}(x)=\cos (n \arccos x), \quad x \in[-1,1],

\]



и рекуррентное соотношение



\[

T_{n+1}(x)=2 x T_{n}(x)-T_{n-1}(x),

\]



с помощью к-рых находят последовательно



\[

\begin{gathered}

T_{0}(x)=1, \quad T_{1}(x)=x, \quad T_{2}(x)=2 x^{2}-1, \\

T_{3}(x)=4 x^{3}-3 x, \quad T_{4}(x)=8 x^{4}-8 x^{2}+1, \\

T_{5}(x)=16 . x^{5}-20 x^{3}+5 x, \ldots

\end{gathered}

\]



Ортонормированные Ч. м.:



\[

\hat{T}_{0}(x)=\frac{1}{V \bar{\pi}} T_{0}(x)=\frac{1}{V \bar{\pi}},

\]



$\hat{T}_{n}(x)=\sqrt{\frac{2}{\pi}} T_{n}(x)=\sqrt{\frac{2}{\pi}} \cos (n \arccos x), n \geqslant 1$.





\end{document}